\section{Conclusion}
\label{sec:conclusion}

We present the first systematic study of how LLMs represent and process French vigesimal numbers.
Our behavioral evaluation shows that \gptmodel has fully mastered the French counting system, achieving 100\% accuracy on conversion, arithmetic, and counting tasks across the vigesimal range.
Our representational analysis of \mistral reveals that this mastery coexists with a distinctive internal organization: vigesimal numbers cluster by linguistic base rather than numeric proximity, with 70\% of nearest neighbors sharing a prefix.
Belgian French's regular decimal system produces representations with superior ordinal structure ($\rho = 0.591$ vs.\ $0.520$), and French number words are more linearly decodable than digit strings ($R^2 = 0.979$ vs.\ $0.943$).

These findings carry two broader implications.
First, behavioral benchmarks alone are insufficient for understanding how models process numbers---probing reveals structure that perfect accuracy masks.
Second, linguistic regularity directly shapes the quality of internal numeric representations, suggesting that the design of a language's number system has consequences for how effectively LLMs encode numerical information.

Future work should extend this analysis to other models, longer number ranges, and causal interventions that test whether the vigesimal representational structure affects downstream reasoning.
Investigating other languages with irregular number systems---such as Danish (which uses a different vigesimal system) or Georgian---would further clarify the relationship between linguistic form and numerical representation in LLMs.
